%====================================================================
% Phongsakon Mark Konrad — Curriculum Vitae
% moderncv, classic style, black. Academic-first ordering.
%====================================================================
\documentclass[11pt,a4paper,sans]{moderncv}

\moderncvstyle{classic}
\moderncvcolor{black}

\usepackage[T1]{fontenc}
\usepackage[utf8]{inputenc}
\usepackage[scale=0.9,vscale=0.92]{geometry}
\usepackage{enumitem}
\setlength{\hintscolumnwidth}{2.4cm}
\setlength{\footskip}{4.5pt}
% tighten vertical rhythm to keep the CV compact
\setlength{\parskip}{0pt}

% black, non-boxed links (hyperref is loaded by moderncv at begin-document)
\AtBeginDocument{\hypersetup{colorlinks=true, urlcolor=black, linkcolor=black, citecolor=black}}

% compact lists for the publication groups
\setlist[enumerate]{leftmargin=1.4em, itemsep=2pt, topsep=2pt, parsep=0pt}

%-------------------- Personal data --------------------
\name{Phongsakon Mark}{Konrad}
\title{Curriculum Vitae}
\address{S\o nderborg, Denmark}{}{}
\phone[mobile]{+45~55~24~01~58}
\email{phongsakon@outlook.dk}
\homepage{phomarkon.github.io}
\social[linkedin]{phongsakonmarkkonrad}
\social[github]{phomarkon}
\extrainfo{\protect\href{https://scholar.google.com/citations?user=zPHjLREAAAAJ}{Google Scholar}}

\begin{document}
\makecvtitle
\vspace{-6mm}

%==================== Summary ====================
\section{Summary}
\cvitem{}{I want to help reach genuine machine intelligence: systems that understand the world and derive its mechanisms the way science let us, not systems that imitate the surface of understanding. I think the dominant paradigm is on the wrong track, and I find scaling laws ugly. My work asks one question, whether the learning process that scaling relies on actually installs internal representations capable of genuine counterfactual reasoning, or only representations that imitate it. I study this through counterfactual reasoning and the geometry and identifiability of learned representations; causal representation learning is home territory. I work phenomenon-first and run the full research lifecycle myself: problem formulation, experimental design, pipeline engineering, manuscript writing, and venue selection. I am a final-year BSc Software Engineering student at SDU, admitted to the MPhil in Machine Learning and Machine Intelligence at the University of Cambridge for 2026, and I bring an unconventional background (four years of military service, three startups, a semester at HKUST) that shapes how I pick problems worth solving.}

%==================== Education ====================
\section{Education}
\cventry{2026--}{MPhil, Machine Learning and Machine Intelligence}{University of Cambridge}{Cambridge, UK}{\textit{admitted; begins Oct 2026}}{}
\cventry{2025}{Exchange Semester, Computer Science \& Engineering}{HKUST}{Hong Kong SAR}{}{Specialisation: COMP4211 Machine Learning, COMP4471 Deep Learning in Computer Vision, COMP4901B Large Language Models, COMP4901Z Reinforcement Learning; COMP6411D Data Visualisation (postgraduate).}
\cventry{2023--2026}{BSc, Software Engineering}{University of Southern Denmark}{S\o nderborg, DK}{GPA $\approx$ 10.6/12 ($\approx$ 3.7/4.0)}{Expected June 2026. BSc thesis (with T.~L.~Adam): \textit{The Trader's Trinity: Forecasting Models, RL Agents, and LLM Judges for Day-Ahead Markets}.}

%==================== Selected Publications ====================
\section{Selected Publications}
\cvitem{2026}{\textbf{P.~M.~Konrad}, T.~Tanyel, S.~Ayvaz. \emph{Self-Reports Do Not Identify Self-Models: An Identifiability Test for Counterfactual Reports}. Accepted (\textbf{oral}), \textit{PhilML Workshop, ICML 2026}. \href{https://phomarkon.github.io/papers/ICML_2026_Workshop_Counterfactual_Self_Reports_Not_Well_Posed.pdf}{[PDF]}\newline {\small\textsf{TL;DR.} Self-reports are evidence about behaviour in a prompt, not proof of a self-model. Wrong-source demonstrations pull affect-state reports toward the source family unless mechanism binding holds.}}
\cvitem{2026}{\textbf{P.~M.~Konrad}, T.~Tanyel, S.~Ayvaz. \emph{Decoded but Unused: Instruction Tuning Routes Moral Framing into the Judgment Readout}. Accepted (poster), \textit{ICML 2026 Workshop on Mechanistic Interpretability}. \href{https://phomarkon.github.io/papers/ICML_2026_Decoded_but_Unused.pdf}{[PDF]}\newline {\small\textsf{TL;DR.} Moral framing is linearly decodable in pretrained models but only causally routed to judgment after instruction tuning.}}
\cvitem{2026}{\textbf{P.~M.~Konrad}, T.~Tanyel, S.~Ayvaz. \emph{A Path Already Walked: On Inheriting Network-Neuroscience Tools for Mechanistic Interpretability}. Accepted (virtual poster), \textit{ICML 2026 Workshop on Mechanistic Interpretability}. \href{https://phomarkon.github.io/papers/ICML_2026_Workshop_A_Path_Already_Walked_Network_Neuroscience.pdf}{[PDF]}\newline {\small\textsf{TL;DR.} Mech-interp should import the graph vocabulary of network neuroscience (modularity, rich clubs, motifs) under explicit contracts.}}
\cvitem{2026}{\textbf{P.~M.~Konrad}, T.~Tanyel, S.~Ayvaz. \emph{Routing Subspaces: Auditing Evaluation-to-Deployment Mismatch in Fine-Tuned Language Models}. Under review. \href{https://phomarkon.github.io/papers/NeurIPS_2026_Routing_Subspaces.pdf}{[PDF]}\newline {\small\textsf{TL;DR.} Path-patching localises an eval-vs-deployment signal to a narrow mid-depth attention window. Clamping a low-dimensional subspace closes the gap in 11 of 12 cells.}}
\cvitem{2026}{\textbf{P.~M.~Konrad}, T.~Tanyel, S.~Ayvaz. \emph{Acceptance Cards: A Four-Diagnostic Standard for Safe Fine-Tuning Defense Claims}. Under review. \href{https://arxiv.org/abs/2605.10575}{arXiv:2605.10575}\newline {\small\textsf{TL;DR.} A four-gate evidential standard for safe-fine-tuning defense claims. SafeLoRA fails the full-card pass under matched-action audit.}}

%==================== Research Experience ====================
\section{Research Experience}
\cventry{Jan 2026--}{Research Collaborator}{HKUST DataVISards group}{Hong Kong SAR}{}{Machine learning and data-visualisation research.}
\cventry{Jan 2026--}{Research Collaborator}{SDU Data \& Intelligence Lab}{S\o nderborg, DK}{}{Independently initiate and drive ML research end-to-end: problem formulation, experimental design, pipeline development, manuscript writing, and venue selection. Co-author with Assoc.\ Prof.\ Serkan Ayvaz across interpretability, NLP, medical imaging, and remote sensing.}
\cventry{Sep 2024--Dec 2025}{Research Assistant}{SDU Data \& Intelligence Lab}{S\o nderborg, DK}{}{Built end-to-end ML pipelines; contributed to study design, literature reviews, model development, and manuscript preparation; supported grant proposals and multimodal dataset management.}

%==================== Full Publication List ====================
\section{Full Publication List}
\cvitem{}{\small \textbf{P.~M.~Konrad} is first author except where noted. Workshop, conference, and under-review work are listed separately; nothing is double-counted.}

\subsection{Peer-reviewed and accepted}
\nopagebreak
{\itshape Interpretability and safety (workshop)}
\begin{enumerate}
\item \textbf{P.~M.~Konrad}, T.~Tanyel, S.~Ayvaz (2026). \emph{Self-Reports Do Not Identify Self-Models: An Identifiability Test for Counterfactual Reports}. Accepted (\textbf{oral}), \textit{PhilML Workshop, ICML 2026}. \href{https://phomarkon.github.io/papers/ICML_2026_Workshop_Counterfactual_Self_Reports_Not_Well_Posed.pdf}{[PDF]}
\item \textbf{P.~M.~Konrad}, T.~Tanyel, S.~Ayvaz (2026). \emph{Decoded but Unused: Instruction Tuning Routes Moral Framing into the Judgment Readout}. Accepted (poster), \textit{ICML 2026 Workshop on Mechanistic Interpretability}. \href{https://phomarkon.github.io/papers/ICML_2026_Decoded_but_Unused.pdf}{[PDF]}
\item \textbf{P.~M.~Konrad}, T.~Tanyel, S.~Ayvaz (2026). \emph{A Path Already Walked: On Inheriting Network-Neuroscience Tools for Mechanistic Interpretability}. Accepted (virtual poster), \textit{ICML 2026 Workshop on Mechanistic Interpretability}. \href{https://phomarkon.github.io/papers/ICML_2026_Workshop_A_Path_Already_Walked_Network_Neuroscience.pdf}{[PDF]}
\end{enumerate}
{\itshape AI for software engineering (workshop)}
\begin{enumerate}[start=4]
\item \textbf{P.~M.~Konrad}, T.~L.~Adam, R.~Terrenzi, S.~Ayvaz (2026). \emph{Architecture Without Architects: How AI Coding Agents Shape Software Architecture}. Accepted, SAGAI Workshop, IEEE ICSA 2026. \href{https://arxiv.org/abs/2604.04990}{arXiv:2604.04990}
\item T.~L.~Adam, \textbf{P.~M.~Konrad}, R.~Terrenzi, F.~G.~Lukas, R.~Yilmaz, K.~Sierszecki, S.~Ayvaz (2026). \emph{CAKE: Cloud Architecture Knowledge Evaluation of Large Language Models}. Accepted, KDA-AI Workshop, IEEE ICSA 2026. \href{https://arxiv.org/abs/2604.05755}{arXiv:2604.05755}
\item R.~Terrenzi, \textbf{P.~M.~Konrad}, T.~L.~Adam, S.~Ayvaz (2026). \emph{A Reference Architecture for Agentic Hybrid Retrieval in Dataset Search}. Accepted, SAML Workshop, IEEE ICSA 2026. \href{https://arxiv.org/abs/2604.16394}{arXiv:2604.16394}
\end{enumerate}
{\itshape Conference proceedings}
\begin{enumerate}[start=7]
\item N.-T.~Thinh, Y.~Du, \textbf{P.~M.~Konrad}, A.~Narechania (2026). \emph{Fact-check Your Information (FYI): A Design Probe to Understand How People Actually Fact-check Data-Driven Articles}. Accepted, \textit{IEEE VIS 2026}. \href{https://phomarkon.github.io/papers/VIS_2026_Fact_check_Your_Information_FYI.pdf}{[PDF]}
\item \textbf{P.~M.~Konrad}, T.~Tanyel, S.~Ayvaz (2025). \emph{Beyond Major Floods: Deep Learning for Detecting Shallow Water Inundation in Agricultural Areas}. \textit{KES 2025}, Procedia Computer Science, 270, 301--310. \href{https://www.sciencedirect.com/science/article/pii/S1877050925028194}{[open access]}
\end{enumerate}

\subsection{Preprints and under review}
\nopagebreak
{\itshape Interpretability and safety}
\begin{enumerate}[start=9]
\item \textbf{P.~M.~Konrad}, T.~Tanyel, S.~Ayvaz (2026). \emph{Routing Subspaces: Auditing Evaluation-to-Deployment Mismatch in Fine-Tuned Language Models}. Under review. \href{https://phomarkon.github.io/papers/NeurIPS_2026_Routing_Subspaces.pdf}{[PDF]}
\item \textbf{P.~M.~Konrad}, T.~Tanyel, S.~Ayvaz (2026). \emph{Acceptance Cards: A Four-Diagnostic Standard for Safe Fine-Tuning Defense Claims}. Under review. \href{https://arxiv.org/abs/2605.10575}{arXiv:2605.10575}
\item \textbf{P.~M.~Konrad}, T.~L.~Adam, A.~C.~H.~Merrild, R.~de~Rosa, R.~Terrenzi, T.~Tanyel, S.~Ayvaz (2026). \emph{The Open-Box Fallacy: Why AI Deployment Needs a Calibrated Verification Regime}. Under review. \href{https://arxiv.org/abs/2605.10601}{arXiv:2605.10601}
\item \textbf{P.~M.~Konrad}, Y.~Du, S.~Ayvaz (2026). \emph{Grounded Auditing as an Evaluation Policy: A Matched-Action Protocol and Stress Test}. Under review. \href{https://phomarkon.github.io/papers/NeurIPS_2026_Grounded_Auditing_Matched_Action_Protocol.pdf}{[PDF]}
\item \textbf{P.~M.~Konrad}, S.~Ayvaz (2026). \emph{When Does Chain-of-Thought Improve Safety? Evidence from 18 Models across 5 Families}. Under review. \href{https://phomarkon.github.io/papers/COLM_2026_When_Does_Chain_of_Thought_Improve_Safety.pdf}{[PDF]}
\item \textbf{P.~M.~Konrad}, Y.~Du, T.~Tanyel, S.~Ayvaz (2026). \emph{How Much of a Probe Direction Lives in the Input Embedding? An Audit of Four Published Claims}. Under review. \href{https://phomarkon.github.io/papers/EMNLP_2026_Probe_Direction_Input_Embedding_Audit.pdf}{[PDF]}
\item \textbf{P.~M.~Konrad}, Y.~Du, T.~Tanyel, S.~Ayvaz (2026). \emph{Abstention and Refusal Are Separable Directions in the Residual Stream}. Under review. \href{https://phomarkon.github.io/papers/EMNLP_2026_Abstention_Refusal_Separable_Directions.pdf}{[PDF]}
\end{enumerate}
{\itshape Applied ML and other}
\begin{enumerate}[start=15]
\item \textbf{P.~M.~Konrad}, A.-A.~Popa, Y.~Sabzehmeidani, L.~Zhong, M.~Tripathy, A.~Constantinescu, E.~A.~Liehn, S.~Ayvaz (2025). \emph{Challenges in Deep Learning-Based Small Organ Segmentation: A Benchmarking Perspective for Medical Research with Limited Datasets}. Under revision, \textit{Biomedical Signal Processing and Control}. \href{https://arxiv.org/abs/2509.05892}{arXiv:2509.05892}
\item \textbf{P.~M.~Konrad}, C.~H.~Kunstmann-Olsen, J.~Fiutowski, S.~Ayvaz (2025). \emph{Non-Destructive Prediction of Fruit Ripeness and Firmness Using Hyperspectral Imaging and Lightweight Machine Learning Models}. Under review. \href{https://arxiv.org/abs/2604.22788}{arXiv:2604.22788}
\end{enumerate}

%==================== Teaching ====================
\section{Teaching}
\cventry{2026}{Teaching Assistant, Artificial Intelligence (BSc)}{University of Southern Denmark}{}{}{Exercise sessions and student supervision; designed hands-on lab assignments; built an \href{https://phomarkon.github.io/ai101/}{interactive AI course} (step-by-step animations) used as supplementary material.}

%==================== Academic Service ====================
\section{Academic Service}
\cvitem{2025--}{\textbf{Student Member}, Educational Committee, Software Engineering Programme, SDU. Represent students in curriculum and quality-assurance discussions on course content, assessment, and study environment.}

%==================== Honors and Awards ====================
\section{Honors and Awards}
\cvitem{2026}{\textbf{William Demant Fonden Grant} supporting graduate study.}
\cvitem{2024, 2025}{\textbf{Top-10}, Danish National Championship in AI (DM i AI), competing solo against 50+ teams.}
\cvitem{2023}{\textbf{1st place}, SDU Case Competition (48-hour sustainability challenge, Danfoss and Linak cases).}
\cvitem{2021}{\textbf{Commendation and performance bonus}, Bundeswehr, for exemplary and independent service.}

%==================== Professional Experience ====================
\section{Professional Experience}
\cventry{2026--}{Founder}{SaturoLabs}{Denmark}{}{Products at the intersection of AI and software engineering. Live: \href{https://claudeboyz.com}{claudeboyz.com}, \href{https://getproofz.com}{getproofz.com}, \href{https://dreambear.app}{dreambear.app}.}
\cventry{2024}{CTO \& Co-Founder}{Tutora ApS}{S\o nderborg, DK}{}{Led end-to-end development of the company's web application.}
\cventry{2022--2024}{Co-CEO \& Co-Founder}{Yeager GmbH}{Remote}{}{Built \emph{stabil.ai}, an AI powerlifting-training app; personalised plans via MRV/MEV with real-time feedback.}
\cventry{2017--2021}{Staff Duty Soldier}{Bundeswehr (German Armed Forces)}{Gl\"ucksburg, DE}{}{Led a small HR team administering 600+ soldiers; twice decorated.}

%==================== Skills ====================
\section{Technical Skills}
\cvitem{Languages}{Python, JavaScript, SQL}
\cvitem{ML}{PyTorch, Hugging Face, TransformerLens, scikit-learn, pandas, NumPy; Optuna, Weights \& Biases}
\cvitem{Tooling}{Git, Docker, Google Cloud Platform; React, React Native, Node.js}

%==================== Languages ====================
\section{Languages}
\cvitem{Fluent}{English (professional), German (native), Thai (native)}
\cvitem{Basic}{Danish}

\end{document}
